\section{Conclusion} \label{sec:conclusion}

This paper documents the calibration-production algorithms, instrument signature removal (ISR) workflow, and quality-assurance framework used to generate DP2 calibration products for LSSTCam.
DP2 is the first end-to-end demonstration of data processing with on-sky LSSTCam dataset, and it establishes a reproducible reference for the first stage of processing in Data Release Processing (DRP) pipeline.

We organize calibration production and ISR around the electronic detector model introduced in \S\ref{sec:lsstcam:detector-model}, where each instrument artifact get its own associated calibration dataset and correction algorithm. Calibrations are built on top of other calibrations (summarized in Figure~\ref{fig:calib-dependency}), requiring intermidate applications of ISR (summarized in Figure~\ref{fig:isr-pipeline}) .
In general, the pipelines are derived and applied in the reverse chronological order of the photon transfer chain.
For DP2, we present the science ISR pipeline (Figure~\ref{fig:isr-pipeline}, \S\ref{sec:isr-science}) and the production algorithms for crosstalk, defect masking, non-linearity, deferred charge (CTI), the photon transfer curve (PTC), the BFE, and the classical bias, dark, and flat-field calibrations (\S\ref{sec:calibration-production}).

Where detailed validation is already in place, DP2 calibrations meet or approach the requirements needed for commissioning-era science.
These products are accompanied by a verification, acceptance, and certification workflow (\S\ref{sec:verify-accept-certify}) intended to make calibration quality auditable, and we record the relevant provenance in exposure dataset metadata for the science community (\S\ref{sec:isr-science:provenance}).

DP2 nevertheless reflects commissioning-era conditions: calibration data were collected before dome completion and full thermal stabilization (\S\ref{sec:intro}), detector operating modes and sequencer configurations evolved across epochs (\S\ref{sec:discussion}), and calibrations and correction tools for some artifacts remain under active development (\S\ref{sec:limitations}).
Users should treat residual focal-plane non-uniformities, epoch-dependent systematics, and incomplete coverage effects as known limitations when interpreting DP2 data products and deciding paths/configurations for science analyses.

Looking toward Data Release~1 (DR1), we expect iterative improvements in three areas: \emph{algorithms} (refined CTI and trap modeling, expanded color-dependent BFE treatment), \emph{operations} (more stable dome and calibration-system performance, clearer calibration-epoch boundaries, and expanded automation through verification and acceptance pipelines), and \emph{diagnostics} (trending analyses, regression and alerts for calibration changes against baselines, and quality metrics on science-images for downstream pipeline acceptance criteria).
The DP2 workflow documented here is the baseline from which those improvements will be measured; maintaining versioned, reproducible calibration production will remain essential as LSSTCam transitions from data previews during commissioning to data releases over the full 10-year survey.

\vskip 0.4in
